\chapter{Referencial Teórico}

\section{Ensino de Computação, inteligência artificial e formação na Educação Básica}

O ensino de Computação na Educação Básica não se limita ao uso de ferramentas digitais. Ele envolve conhecimentos e práticas que permitem compreender o funcionamento de sistemas computacionais. Recursos cotidianos baseados em inteligência artificial aproximam essa discussão da experiência dos estudantes, mas essa proximidade não substitui o ensino dos princípios que orientam esses recursos.

A Resolução CNE/CEB nº 1/2022 estabelece normas para a Computação como complemento à BNCC e atribui aos sistemas de ensino a definição de parâmetros e abordagens para sua implementação \cite{brasil2022computacao}.

O ensino de Computação difere da informática na educação, em que o computador apoia a aprendizagem de conteúdos de diversas áreas e o foco recai sobre sua inserção pedagógica \cite{valente2002uso}. A plataforma deste projeto também será um recurso didático, mas o objeto central de aprendizagem não é operar a interface nem estudar um conteúdo externo com auxílio do computador.

A proposta focaliza conceitos e práticas da Computação relacionados a dados, características, modelos, regras, parâmetros, treinamento, previsão, classificação, agrupamento e limites de sistemas de inteligência artificial. Espera-se que os estudantes interpretem como dados e regras participam da produção de uma saída, comparem finalidades de diferentes modelos, antecipem efeitos de modificações e expliquem os resultados observados.

A Matemática apoia essa leitura por meio de gráficos, funções, distância, probabilidade, média e erro, sem substituir o conteúdo computacional que organiza as atividades. A Computação na Educação Básica se estrutura pelos eixos Pensamento Computacional, Mundo Digital e Cultura Digital, para os quais as Diretrizes da Sociedade Brasileira de Computação e a Resolução CNE/CEB nº 1/2022 oferecem referências \cite{sbc2018diretrizes,brasil2022computacao}. Os eixos não serão tratados como conteúdos separados, mas evitam que a plataforma se reduza a controles técnicos. No Ensino Médio, as habilidades da norma ligadas à inteligência artificial, aos modelos computacionais e à análise de dados orientam a seleção dos conteúdos \cite{brasil2022computacao}.

O pensamento computacional participa da proposta como forma de organizar e analisar problemas, dados, regras e procedimentos, e não como sinônimo de qualquer atividade realizada com computador. \citeonline{wing2006computational} chama a atenção para práticas como abstração e decomposição na formulação de soluções computacionais. Neste estudo, essa aproximação aparece quando o estudante separa uma situação em dados disponíveis, saída desejada, regra ou modelo e resultado produzido.

O estudo não tratará essa mobilização como medida de um desenvolvimento amplo de pensamento computacional nem como base para afirmar melhoria em outros componentes curriculares. Ela ajuda a interpretar o conteúdo específico da sequência.

No Mundo Digital, os estudantes analisam como modelos processam representações numéricas e executam regras para produzir respostas, percebendo que uma decisão automatizada depende de elementos que podem ser analisados. A Cultura Digital aparece no exame do uso de sistemas de inteligência artificial presentes no cotidiano.

\citeonline{touretzky2019envisioning} defendem que a educação em inteligência artificial para a Educação Básica deve ajudar os estudantes a compreender ideias que orientam sistemas inteligentes e seus impactos. O documento da AI4K12 sobre \textit{machine learning} \cite{touretzky2023machine} situa a aprendizagem a partir de dados como uma das grandes ideias da inteligência artificial e inclui, para a faixa do Ensino Médio, a diferenciação entre aprendizagem supervisionada e não supervisionada. Essas referências não são um currículo pronto para a realidade brasileira, mas oferecem linguagem conceitual para selecionar conhecimentos que não dependam de programação avançada.

\citeonline{camada2020ensino} identificaram possibilidades de pesquisa e aplicação para o ensino de inteligência artificial na Educação Básica a partir de uma revisão sistemática. \citeonline{duraes2021ensino} relatam uma sequência didática sobre fundamentos de IA com estudantes do Ensino Médio integrado. Os estudos mostram que o tema já vem sendo trabalhado e reforçam a necessidade de considerar o tempo escolar, o repertório matemático, a mediação docente e a infraestrutura disponível.

Neste projeto, a inteligência artificial permite discutir como sistemas computacionais dependem de decisões humanas de modelagem. Compreender de onde vêm os dados e o que muda quando um parâmetro é alterado permite ao estudante questionar o resultado que recebe, em vez de apenas aceitá-lo. A questão educacional não será se o estudante consegue usar uma ferramenta generativa para terminar uma tarefa escolar. Será verificar se consegue explicar, em situações delimitadas, como alterações nos dados, nos parâmetros ou nas regras participam de previsões, classificações e agrupamentos, o que mantém a proposta vinculada ao ensino de Computação.

\section{\textit{Machine learning} introdutório no Ensino Médio}

O \textit{machine learning} ocupa lugar específico entre os conteúdos de inteligência artificial porque permite examinar como computadores identificam padrões em dados e os utilizam para produzir respostas em tarefas definidas. Ao mapearem unidades instrucionais destinadas à escola, \citeonline{marques2020teaching} observaram que vários materiais apresentam apenas aspectos básicos de \textit{machine learning} ou deixam partes do funcionamento do modelo em nível abstrato.

A questão pedagógica não é expor cada algoritmo em todos os detalhes técnicos. O estudante precisa compreender o suficiente para não tratar o resultado do modelo como uma resposta sem história.

\citeonline{martins2023findings} reúnem evidências de que estudantes iniciantes do Ensino Médio podem compreender e aplicar conceitos, algoritmos e tarefas básicas de \textit{machine learning}. \citeonline{sanusi2023systematic}, em revisão sobre a Educação Básica, também identificam experiências em diferentes níveis escolares, com presença expressiva do Ensino Médio. As pesquisas indicam lacunas na integração curricular, na formação de professores e na avaliação das intervenções. A adequação à faixa etária não depende apenas da idade dos participantes, mas do que se escolhe ensinar, de como as relações são representadas e da ajuda oferecida durante a atividade.

Esses estudos não permitem tratar o Ensino Médio como inadequado, de partida, para discutir aprendizagem a partir de dados e modelos simples. Também não permitem supor que qualquer atividade digital sobre inteligência artificial produzirá compreensão. As intervenções analisadas têm recursos, durações e objetivos distintos. Uma experiência conduzida em uma escola específica não se converte automaticamente em receita para outra. A literatura permite formular uma proposta viável, mas não antecipa seus resultados.

A sequência não deve começar pela expressão matemática de um algoritmo, pois isso pode afastar estudantes que ainda consolidam a leitura de gráficos e noções de proporcionalidade. A manipulação de controles e a repetição de nomes técnicos também podem ocultar o raciocínio que se pretende desenvolver. A introdução ao tema precisa partir de problemas reconhecíveis, avançar para dados e representações examináveis e aproximar, em seguida, a regra que gera uma saída. Essa progressão cria uma porta de entrada compatível com o nível de escolaridade previsto, sem diminuir o conteúdo de Computação.

O vocabulário técnico não é a única dificuldade esperada. O estudante pode precisar ler dois eixos, localizar um ponto, comparar distâncias ou reconhecer o que uma cor, uma curva ou uma regra representa. A leitura da visualização precisa ser ensinada como parte da atividade, pois observar uma mudança não significa explicar a relação que a produziu.

Há também um obstáculo matemático. Noções de função, média, probabilidade e distância nem sempre estão estabilizadas no repertório da turma. A sequência não resolverá, em oito encontros, lacunas acumuladas na escolarização matemática. Ela retomará essas noções quando forem necessárias para responder a uma pergunta concreta sobre o modelo, sem ocultar a Matemática nem antecipar uma formalização que torne a atividade inacessível.

Essa retomada não transforma o percurso em revisão extensa de Matemática antes do conteúdo de Computação. A necessidade aparece dentro de cada problema, como na comparação entre previsão e valor observado na regressão linear ou na distância entre exemplos no \textit{k}-NN. A mediação deverá identificar quando uma retomada curta basta e quando a dificuldade impede que a turma prossiga com a interpretação planejada. A seção 3.3 detalha como a Matemática entra nas representações.

Uma interface visual pode reduzir algumas barreiras e, ao mesmo tempo, ocultar outras. Um controle que altera o número de vizinhos torna uma mudança visível, mas não explica, por si só, por que certos pontos passam a decidir uma categoria. O recurso visual precisa ser acompanhado por mediação orientada, como se discute na seção 3.3.

A plataforma trabalha regressão linear, regressão logística, \textit{k}-vizinhos mais próximos, árvore de decisão e \textit{k}-means. A seleção permite discutir previsão de valores numéricos, classificação e formação de grupos sem rótulos prévios. Em cada módulo, o estudante deverá reconhecer o problema tratado, levantar uma hipótese sobre o efeito de uma alteração e explicar, após o experimento, a relação entre dados, parâmetro, gráfico e resultado.

Os modelos foram escolhidos porque mostram mecanismos distintos em representações examináveis por estudantes iniciantes. A regressão linear permite discutir tendência, previsão, resíduo e erro. A regressão logística separa probabilidade e classificação, pois uma variação no limiar altera a decisão sem modificar o valor calculado. O \textit{k}-NN aproxima a classificação pelo confronto com exemplos semelhantes. A árvore de decisão torna visível uma sequência de condições. O \textit{k}-means permite discutir grupos quando não existe uma categoria indicada para cada ponto.

O recorte aproxima duas formas de aprendizagem a partir de dados. Em regressão linear, regressão logística, \textit{k}-NN e árvore de decisão, os exemplos trazem uma resposta conhecida que orienta a tarefa de prever um valor ou escolher uma categoria. No \textit{k}-means, não há categoria fornecida previamente, e o objetivo é formar grupos a partir de uma medida de proximidade. \citeonline{touretzky2023machine} situam essa diferenciação entre aprendizagem supervisionada e não supervisionada como uma das ideias centrais para a educação em IA. Neste projeto, ela será abordada com exemplos sintéticos e visualizações.

Os cinco módulos não devem ser assimilados com a mesma velocidade ou pelo mesmo caminho. A regressão linear pode se aproximar de conteúdos matemáticos já conhecidos, enquanto o agrupamento sem rótulos pode exigir mais tempo para que os estudantes diferenciem descobrir grupos e separar categorias previamente definidas. A retomada no encontro seguinte e a comparação coletiva no sétimo encontro evitam tratar cada modelo como episódio isolado.

O nível introdutório não exige ocultar os limites dos modelos. Uma reta resume a tendência apenas dentro do conjunto observado e pode ser afetada por pontos atípicos. Uma porcentagem representada na regressão logística não é certeza.

No \textit{k}-NN, a decisão depende dos exemplos guardados, da escala das características e do número de vizinhos escolhido. Em árvores, regras mais detalhadas não garantem melhor comportamento em casos novos. No \textit{k}-means, os grupos dependem do número de centroides e de sua posição inicial. Esses limites serão tratados como questões de interpretação. O objetivo é evitar a promessa de precisão automática e uma abordagem moralizante e abstrata sobre inteligência artificial. Quando o estudante altera um dado, um limiar, uma condição ou um centroide, deve poder observar uma consequência e discutir de que regra ela decorre.

\section{Representações, visualização e experimentação guiada}

Ferramentas visuais permitem introduzir \textit{machine learning} a estudantes iniciantes. \citeonline{gressevonwangenheim2021visual} identificaram ferramentas que favorecem a construção interativa de modelos, embora várias se concentrem na classificação de imagens e em atividades curtas. O produto proposto trabalha com dados tabulares, coletados pela turma e sintéticos, e inclui regressão, classificação e agrupamento. Essa diferença de escopo não resolve, por si só, as limitações apontadas pela literatura.

A interação só ganha sentido pedagógico quando o estudante investiga uma relação. A visualização pode tornar visível algo que permaneceria abstrato em uma explicação verbal, mas também pode dar a impressão de que observar uma mudança basta para compreendê-la. Os elementos visuais precisam ter significado explícito, e a regra do modelo deve estar acessível.

Os dados sintéticos permitem controlar o que será comparado e produzir situações em que a modificação de um elemento gera uma consequência perceptível. Não servem como previsão confiável sobre pessoas reais. Os dados coletados pela turma, nos encontros 2 e 8, mostram de onde vêm os números que serão interpretados. O cenário didático permite perguntar o que mudou, por que mudou e qual elemento do modelo participou da mudança.

Tabela, gráfico, fórmula, texto explicativo e controle interativo mobilizam múltiplas representações. \citeonline{ainsworth2006deft} argumenta que representações diferentes podem cumprir funções complementares durante a aprendizagem, mas exigem tarefas de interpretação e integração. A plataforma não deve apenas reunir recursos na mesma tela. Cada representação precisa responder à pergunta examinada pela turma.

A alteração na inclinação da reta, no limiar de decisão, no número de vizinhos, na condição da árvore ou na posição dos centroides deve vir acompanhada de uma questão que leve o estudante a antecipar e interpretar a mudança. Na regressão linear, a tabela mostra pares de valores, o plano cartesiano apresenta a distribuição dos pontos e a reta permite examinar tendência ou previsão. Na regressão logística, a curva e a probabilidade precisam ser lidas junto ao limiar que transforma uma estimativa em classe. No \textit{k}-NN, o estudante relacionará a posição do ponto consultado, a distância até os exemplos e a votação dos vizinhos. Em árvores, o caminho de regras deve ser associado às condições que o produziram. Nos agrupamentos, os centroides precisam ser entendidos como posições que reorganizam os grupos, e não como marcadores decorativos. A visualização sustenta, em cada caso, a relação entre dado, regra e saída.

O ciclo prever, experimentar, observar e explicar organiza a atividade. Antes de modificar um controle, o estudante formula uma hipótese, por escrito ou oralmente. Durante o experimento, observa o que muda no gráfico, na decisão ou na medida calculada. Depois, explica por que o resultado confirmou ou contrariou sua expectativa. O ciclo evita a tentativa e erro em que parâmetros são movimentados sem uma relação clara entre entrada, regra e saída.

A previsão torna visível uma ideia que poderia permanecer implícita. Se um estudante afirma que aumentar o número de vizinhos não alterará uma classificação e encontra decisão diferente, o professor pode perguntar quais exemplos passaram a ser considerados e por que a votação se alterou. Se prevê que mover um centroide não terá efeito e observa a reorganização dos pontos, a explicação pode retomar a noção de proximidade. A hipótese inicial não precisa estar correta para ter valor pedagógico.

A aplicação ocorrerá no laboratório de informática da escola, com acesso à internet. Os estudantes trabalharão em duplas, cada uma com um computador, e a restrição de acesso fora da escola não impedirá a participação, pois todas as atividades acontecerão nos encontros. Para estudantes com insegurança no manuseio do computador, os controles deverão ser legíveis e indicar o que pode ser modificado em cada momento. Para aqueles que operam a interface com facilidade, o desafio estará na explicação, não na velocidade de mover um controle. A mediação docente e o trabalho em duplas ajudam a evitar que a técnica de uso se torne obstáculo ou substitua a questão conceitual.

O controle interativo não será tratado como decoração ou recompensa após uma explicação pronta. Ele permite comparar uma consequência ao manter alguns elementos estáveis e alterar um dado, parâmetro ou regra. A aula guiada orienta, inicialmente, uma alteração por vez, pois mudanças simultâneas dificultam atribuir uma consequência a uma causa percebida.

Depois da aula guiada, a exploração livre poderá ampliar as comparações. O professor deverá retomar a finalidade do experimento e solicitar explicações que não se limitem a descrever elementos apresentados.

As representações matemáticas serão apresentadas em camadas. Uma situação concreta introduz o problema; a tabela ou o gráfico permite localizar dados e relações; um cálculo elementar explicita uma quantidade relevante; a fórmula, quando utilizada, resume a regra discutida. A fórmula da reta, a distância entre pontos, o limiar, uma condição de árvore ou a média que atualiza um centroide devem explicar uma alteração observável, e não funcionar como requisito de memorização desarticulada.

Apresentar a Matemática em camadas não retira precisão dos conceitos. Define o que é necessário para uma compreensão inicial e o momento em que isso passa a fazer sentido. Antes de calcular a distância entre dois pontos, o estudante precisa reconhecer por que a proximidade importa no \textit{k}-NN. Antes de discutir o erro de uma reta, precisa comparar uma previsão com um valor observado.

Quando a fórmula entra depois dessas experiências, ela deixa de ser informação isolada e passa a nomear uma relação que o estudante já pode examinar no experimento.

\section{Avaliação da compreensão durante a intervenção}

Avaliar o ensino de \textit{machine learning} na Educação Básica permanece uma questão em aberto. \citeonline{rauber2023assessing} levantaram avaliações que vão de questionários a tarefas baseadas em desempenho e destacaram a necessidade de abordagens mais abrangentes.

\citeonline{gressevonwangenheim2022proposal} apresentam uma rubrica baseada em desempenho para conceitos e práticas de \textit{machine learning}. A proposta mostra a necessidade de transformar expectativas de aprendizagem em critérios observáveis. Embora esteja situada em classificação de imagens, a proposta exige adaptação ao recorte desta pesquisa. Reconhecer o nome de um algoritmo não demonstra, por si só, que o estudante compreende a tarefa que ele resolve ou a relação entre dados e saída. Uma resposta numérica correta também pode resultar de tentativa, leitura superficial de um gráfico ou repetição de uma fala do professor. Os instrumentos devem solicitar justificativas breves que permitam examinar como o estudante chega a uma conclusão.

A compreensão não será medida pelo acerto em uma questão de múltipla escolha. O diagnóstico inicial e o desafio final individual oferecerão produções comparáveis. Os registros breves do professor-pesquisador ajudarão a recuperar condições de uso, dúvidas recorrentes e intervenções realizadas durante as aulas.

O diagnóstico inicial não selecionará estudantes nem medirá deficiência. Ele mostrará ideias de partida, por exemplo, quando alguém atribui uma decisão à inteligência artificial sem mencionar dados ou regras, confunde previsão com classificação ou supõe que agrupamentos já possuem categorias definidas. O desafio final retomará dimensões semelhantes em uma situação nova, com uma tabela organizada pelos próprios estudantes, sem depender da repetição literal dos exemplos explorados na plataforma.

\citeonline{black1998assessment} discutem a avaliação formativa como parte do ensino, pois as informações produzidas pelos estudantes podem orientar decisões docentes. As perguntas disponíveis na plataforma terão papel de devolutiva imediata e de organização da conversa em sala. As respostas não serão coletadas digitalmente. A análise se concentrará nos instrumentos individuais e nos registros docentes previstos para a intervenção.

Durante a exploração, o professor poderá usar as previsões para identificar uma dificuldade que exija retomada, sem registrar cada fala como dado de pesquisa. Se muitas duplas interpretarem uma probabilidade de 70\% como garantia de ocorrência, a aula poderá interromper a manipulação para discutir a diferença entre estimativa e certeza. Se houver confusão entre \textit{k}-NN e \textit{k}-means, a comparação poderá recuperar a presença ou ausência de rótulos nos exemplos.

Essas intervenções pertencem ao ensino. O registro posterior apenas contextualiza as produções individuais.

Os critérios organizam a análise em quatro dimensões: reconhecer a origem dos dados e a finalidade da tarefa, antecipar consequências de uma alteração, interpretar representações e discutir limites de uma decisão automatizada. Eles não servirão para ordenar estudantes nem produzir um índice isolado de competência. Orientarão a leitura das explicações, distinguindo respostas que apenas nomeiam elementos, descrevem mudanças ou articulam dados, procedimento e resultado.

Uma situação sobre consumo fictício de energia, por exemplo, poderá exigir que o estudante diferencie temperatura e consumo, em vez de chamar todos os valores simplesmente de dados. Outras situações poderão solicitar a antecipação do efeito de alterar um ponto, um limiar, um número de vizinhos, uma condição ou um centroide. A leitura de tabela, gráfico, caminho de decisão e agrupamento será examinada em relação à explicação apresentada, e não apenas pelo acerto isolado.

As respostas podem revelar níveis diferentes de elaboração sem exigir uma escala rígida de pontuação. Há diferença entre escrever que ``mudou porque o programa mudou'' e explicar que a alteração de limiar passou a classificar como alerta valores antes abaixo da regra estabelecida. Esses contrastes ajudam a ler as explicações, sem produzir rótulos fixos para caracterizar estudantes.

Os resultados serão apresentados em sínteses por dimensão, frequências descritivas quando ajudarem a situar o conjunto e poucos excertos anonimizados que ilustrem explicações contrastantes. As respostas individuais não serão comparadas entre alunos, e a pesquisa não pretende demonstrar uma relação causal isolada entre plataforma e aprendizagem. A discussão poderá tratar da trajetória de uma intervenção específica, das evidências produzidas em seu interior e dos limites que impedem extrapolá-las para outros contextos.

\section{Síntese do referencial e posicionamento da proposta}

A literatura e os documentos curriculares definem o posicionamento da proposta. O \textit{machine learning} será tratado como conteúdo introdutório de Computação, e não como pretexto para ensinar o uso instrumental de aplicativos. A sequência seleciona ideias fundamentais, tarefas e limites observáveis, sem antecipar uma formação universitária em algoritmos. Os dados coletados pela turma e os cenários sintéticos servem ao mesmo propósito, o de permitir que o estudante saiba de onde vêm os números e o que muda quando um parâmetro é alterado. A visualização organizará uma comparação em que o estudante formula uma hipótese, observa uma consequência e explica a relação entre dado, regra e resultado.

A análise buscará evidências em produções individuais contextualizadas. Não se limitará à medida de acertos nem prometerá um efeito causal amplo.

A literatura registra propostas, ferramentas e experiências para o ensino de IA e \textit{machine learning} na Educação Básica, mas ainda aponta a necessidade de avaliações mais rigorosas, clareza sobre conteúdos e estratégias e atenção aos contextos de aplicação (\MakeUppercase{MARTINS; GRESSE VON WANGENHEIM}, \citeyear{martins2023findings}; \MakeUppercase{RAUBER; GRESSE VON WANGENHEIM}, \citeyear{rauber2023assessing}; \MakeUppercase{SANUSI} et al., \citeyear{sanusi2023systematic}). Este projeto não pretende preencher essa lacuna sozinho nem apresentar uma plataforma inédita. A contribuição será tornar examinável a articulação entre cinco modelos, visualizações, sequência didática, instrumentos e critérios de análise, em uma turma de Ensino Médio em que o pesquisador é o professor.

A qualidade da pesquisa dependerá da coerência entre referencial, produto e método. Cada recurso do site deverá corresponder a um objetivo de aprendizagem, cada objetivo deverá aparecer em uma atividade e em um instrumento, e as conclusões deverão respeitar o que as evidências permitirem afirmar.
